\section{Elementary results on triangulated categories}
\label{section-elementary-results}

\noindent
Most of the results in this section are proved for pre-triangulated categories
and a fortiori hold in any triangulated category.

\begin{lemma}
\label{lemma-composition-zero}
Let $\mathcal{D}$ be a pre-triangulated category.
Let $(X, Y, Z, f, g, h)$ be a distinguished triangle.
Then $g \circ f = 0$,
$h \circ g = 0$ and $f[1] \circ h = 0$.
\end{lemma}

\begin{proof}
By TR1 we know $(X, X, 0, 1, 0, 0)$ is a distinguished triangle.
Apply TR3 to
$$
\xymatrix{
X \ar[r] \ar[d]^1 &
X \ar[r] \ar[d]^f &
0 \ar[r] \ar@{-->}[d] &
X[1] \ar[d]^{1[1]} \\
X \ar[r]^f &
Y \ar[r]^g &
Z \ar[r]^h &
X[1]
}
$$
Of course the dotted arrow is the zero map. Hence the commutativity of
the diagram implies that $g \circ f = 0$. For the other cases
rotate the triangle, i.e., apply TR2.
\end{proof}

\begin{lemma}
\label{lemma-representable-homological}
Let $\mathcal{D}$ be a pre-triangulated category.
For any object $W$ of $\mathcal{D}$ the functor
$\Hom_\mathcal{D}(W, -)$ is homological, and the functor
$\Hom_\mathcal{D}(-, W)$ is cohomological.
\end{lemma}

\begin{proof}
Consider a distinguished triangle $(X, Y, Z, f, g, h)$.
We have already seen that $g \circ f = 0$, see
Lemma \ref{lemma-composition-zero}.
Suppose $a : W \to Y$ is a morphism such that $g \circ a = 0$.
Then we get a commutative diagram
$$
\xymatrix{
W \ar[r]_1 \ar@{..>}[d]^b &
W \ar[r] \ar[d]^a &
0 \ar[r] \ar[d]^0 &
W[1] \ar@{..>}[d]^{b[1]} \\
X \ar[r] & Y \ar[r] & Z \ar[r] & X[1]
}
$$
Both rows are distinguished triangles (use TR1 for the top row).
Hence we can fill the dotted arrow $b$ (first rotate using TR2,
then apply TR3, and then rotate back). This proves the lemma.
\end{proof}

\begin{lemma}
\label{lemma-third-isomorphism-triangle}
Let $\mathcal{D}$ be a pre-triangulated category.
Let
$$
(a, b, c) : (X, Y, Z, f, g, h) \to (X', Y', Z', f', g', h')
$$
be a morphism of distinguished triangles. If two among $a, b, c$
are isomorphisms so is the third.
\end{lemma}

\begin{proof}
Assume that $a$ and $c$ are isomorphisms.
For any object $W$ of $\mathcal{D}$ write
$H_W( - ) = \Hom_\mathcal{D}(W, -)$.
Then we get a commutative diagram of abelian groups
$$
\xymatrix{
H_W(Z[-1]) \ar[r] \ar[d] &
H_W(X) \ar[r] \ar[d] &
H_W(Y) \ar[r] \ar[d] &
H_W(Z) \ar[r] \ar[d] &
H_W(X[1]) \ar[d] \\
H_W(Z'[-1]) \ar[r] &
H_W(X') \ar[r] &
H_W(Y') \ar[r] &
H_W(Z') \ar[r] &
H_W(X'[1])
}
$$
By assumption the right two and left two vertical arrows are bijective.
As $H_W$ is homological by
Lemma \ref{lemma-representable-homological}
and the five lemma
(Homology, Lemma \ref{homology-lemma-five-lemma})
it follows that the middle vertical arrow is an isomorphism.
Hence by Yoneda's lemma, see
Categories, Lemma \ref{categories-lemma-yoneda}
we see that $b$ is an isomorphism.
This implies the other cases by rotating (using TR2).
\end{proof}

\begin{remark}
\label{remark-special-triangles}
Let $\mathcal{D}$ be an additive category with translation functors $[n]$
as in Definition \ref{definition-triangle}. Let us call a triangle
$(X, Y, Z, f, g, h)$ {\it special}\footnote{This is nonstandard notation.}
if for every object $W$ of $\mathcal{D}$
the long sequence of abelian groups
$$
\ldots \to
\Hom_\mathcal{D}(W, X) \to
\Hom_\mathcal{D}(W, Y) \to
\Hom_\mathcal{D}(W, Z) \to
\Hom_\mathcal{D}(W, X[1]) \to \ldots
$$
is exact.  The proof of Lemma \ref{lemma-third-isomorphism-triangle}
shows that if
$$
(a, b, c) : (X, Y, Z, f, g, h) \to (X', Y', Z', f', g', h')
$$
is a morphism of special triangles and if two among $a, b, c$
are isomorphisms so is the third. There is a dual statement for
{\it co-special} triangles, i.e., triangles which turn into long
exact sequences on applying the functor $\Hom_\mathcal{D}(-, W)$.
Thus distinguished triangles are special and co-special, but in
general there are many more (co-)special triangles, than there are
distinguished triangles.
\end{remark}

\begin{lemma}
\label{lemma-third-map-square-zero}
Let $\mathcal{D}$ be a pre-triangulated category.
Let
$$
(0, b, 0), (0, b', 0) : (X, Y, Z, f, g, h) \to (X, Y, Z, f, g, h)
$$
be endomorphisms of a distinguished triangle. Then $bb' = 0$.
\end{lemma}

\begin{proof}
Picture
$$
\xymatrix{
X \ar[r] \ar[d]^0 &
Y \ar[r] \ar[d]^{b, b'} \ar@{..>}[ld]^\alpha &
Z \ar[r] \ar[d]^0 \ar@{..>}[ld]^\beta &
X[1] \ar[d]^0 \\
X \ar[r] & Y \ar[r] & Z \ar[r] & X[1]
}
$$
Applying
Lemma \ref{lemma-representable-homological}
we find dotted arrows $\alpha$ and $\beta$ such that
$b' = f \circ \alpha$ and $b = \beta \circ g$. Then
$bb' = \beta \circ g \circ f \circ \alpha = 0$
as $g \circ f = 0$ by
Lemma \ref{lemma-composition-zero}.
\end{proof}

\begin{lemma}
\label{lemma-third-map-idempotent}
Let $\mathcal{D}$ be a pre-triangulated category.
Let $(X, Y, Z, f, g, h)$ be a distinguished triangle.
If
$$
\xymatrix{
Z \ar[r]_h \ar[d]_c & X[1] \ar[d]^{a[1]} \\
Z \ar[r]^h & X[1]
}
$$
is commutative and $a^2 = a$, $c^2 = c$, then there exists a
morphism $b : Y \to Y$ with $b^2 = b$ such that
$(a, b, c)$ is an endomorphism of the triangle $(X, Y, Z, f, g, h)$.
\end{lemma}

\begin{proof}
By TR3 there exists a morphism $b'$ such that
$(a, b', c)$ is an endomorphism of $(X, Y, Z, f, g, h)$.
Then $(0, (b')^2 - b', 0)$ is also an endomorphism. By
Lemma \ref{lemma-third-map-square-zero}
we see that $(b')^2 - b'$ has square zero.
Set $b = b' - (2b' - 1)((b')^2 - b') = 3(b')^2 - 2(b')^3$.
A computation shows that $(a, b, c)$ is an endomorphism and
that $b^2 - b = (4(b')^2 - 4b' - 3)((b')^2 - b')^2 = 0$.
\end{proof}

\begin{lemma}
\label{lemma-cone-triangle-unique-isomorphism}
Let $\mathcal{D}$ be a pre-triangulated category.
Let $f : X \to Y$ be a morphism of $\mathcal{D}$.
There exists a distinguished triangle $(X, Y, Z, f, g, h)$ which
is unique up to (nonunique) isomorphism of triangles.
More precisely, given a second such distinguished triangle
$(X, Y, Z', f, g', h')$ there exists an isomorphism
$$
(1, 1, c) : (X, Y, Z, f, g, h) \longrightarrow (X, Y, Z', f, g', h')
$$
\end{lemma}

\begin{proof}
Existence by TR1. Uniqueness up to isomorphism by TR3 and
Lemma \ref{lemma-third-isomorphism-triangle}.
\end{proof}

\begin{lemma}
\label{lemma-uniqueness-third-arrow}
Let $\mathcal{D}$ be a pre-triangulated category. Let
$$
(a, b, c) : (X, Y, Z, f, g, h) \to (X', Y', Z', f', g', h')
$$
be a morphism of distinguished triangles. If one of the following
conditions holds
\begin{enumerate}
\item $\Hom(Y, X') = 0$,
\item $\Hom(Z, Y') = 0$,
\item $\Hom(X, X') = \Hom(Z, X') = 0$,
\item $\Hom(Z, X') = \Hom(Z, Z') = 0$, or
\item $\Hom(X[1], Z') = \Hom(Z, X') = 0$
\end{enumerate}
then $b$ is the unique morphism from $Y \to Y'$ such that
$(a, b, c)$ is a morphism of triangles.
\end{lemma}

\begin{proof}
If we have a second morphism of triangles $(a, b', c)$
then $(0, b - b', 0)$ is a morphism of triangles. Hence we
have to show: the only morphism $b : Y \to Y'$ such that
$X \to Y \to Y'$ and $Y \to Y' \to Z'$ are zero is $0$.
We will use Lemma \ref{lemma-representable-homological}
without further mention. In particular, condition (3) implies (1).
Given condition (1) if the composition $g' \circ b : Y \to Y' \to Z'$
is zero, then $b$ lifts to a morphism $Y \to X'$ which has to be zero.
This proves (1).

\medskip\noindent
The proof of (2) and (4) are dual to this argument.

\medskip\noindent
Assume (5). Consider the diagram
$$
\xymatrix{
X \ar[r]_f \ar[d]^0 &
Y \ar[r]_g \ar[d]^b &
Z \ar[r]_h \ar[d]^0 \ar@{..>}[ld]^\epsilon &
X[1] \ar[d]^0 \\
X' \ar[r]^{f'} &
Y' \ar[r]^{g'} &
Z' \ar[r]^{h'} &
X'[1]
}
$$
We may choose $\epsilon$ such that $b = \epsilon \circ g$.
Then $g' \circ \epsilon \circ g = 0$ which implies that
$g' \circ \epsilon = \delta \circ h$ for some
$\delta \in \Hom(X[1], Z')$. Since $\Hom(X[1], Z') = 0$
we conclude that $g' \circ \epsilon = 0$. Hence
$\epsilon = f' \circ \gamma$ for some $\gamma \in \Hom(Z, X')$.
Since $\Hom(Z, X') = 0$ we conclude that $\epsilon = 0$
and hence $b = 0$ as desired.
\end{proof}

\begin{lemma}
\label{lemma-third-object-zero}
Let $\mathcal{D}$ be a pre-triangulated category.
Let $f : X \to Y$ be a morphism of $\mathcal{D}$.
The following are equivalent
\begin{enumerate}
\item $f$ is an isomorphism,
\item $(X, Y, 0, f, 0, 0)$ is a distinguished triangle, and
\item for any distinguished triangle $(X, Y, Z, f, g, h)$ we have $Z = 0$.
\end{enumerate}
\end{lemma}

\begin{proof}
By TR1 the triangle $(X, X, 0, 1, 0, 0)$ is distinguished.
Let $(X, Y, Z, f, g, h)$ be a distinguished triangle.
By TR3 there is a map of distinguished triangles
$(1, f, 0) : (X, X, 0) \to (X, Y, Z)$.
If $f$ is an isomorphism, then $(1, f, 0)$ is an isomorphism
of triangles by Lemma \ref{lemma-third-isomorphism-triangle}
and $Z = 0$. Conversely, if $Z = 0$, then $(1, f, 0)$ is an
isomorphism of triangles as well, hence $f$ is an isomorphism.
\end{proof}

\begin{lemma}
\label{lemma-direct-sum-triangles}
Let $\mathcal{D}$ be a pre-triangulated category.
Let $(X, Y, Z, f, g, h)$ and $(X', Y', Z', f', g', h')$ be triangles.
The following are equivalent
\begin{enumerate}
\item $(X \oplus X', Y \oplus Y', Z \oplus Z',
f \oplus f', g \oplus g', h \oplus h')$
is a distinguished triangle,
\item both $(X, Y, Z, f, g, h)$ and $(X', Y', Z', f', g', h')$ are
distinguished triangles.
\end{enumerate}
\end{lemma}

\begin{proof}
Assume (2). By TR1 we may choose a distinguished triangle
$(X \oplus X', Y \oplus Y', Q, f \oplus f', g'', h'')$.
By TR3 we can find morphisms of distinguished triangles
$(X, Y, Z, f, g, h) \to
(X \oplus X', Y \oplus Y', Q, f \oplus f', g'', h'')$
and
$(X', Y', Z', f', g', h') \to
(X \oplus X', Y \oplus Y', Q, f \oplus f', g'', h'')$.
Taking the direct sum of these morphisms
we obtain a morphism of triangles
$$
\xymatrix{
(X \oplus X', Y \oplus Y', Z \oplus Z',
f \oplus f', g \oplus g', h \oplus h')
\ar[d]^{(1, 1, c)} \\
(X \oplus X', Y \oplus Y', Q, f \oplus f', g'', h'').
}
$$
In the terminology of Remark \ref{remark-special-triangles}
this is a map of special triangles (because a direct sum of special
triangles is special) and we conclude
that $c$ is an isomorphism. Thus (1) holds.

\medskip\noindent
Assume (1). We will show that $(X, Y, Z, f, g, h)$ is a distinguished
triangle. First observe that $(X, Y, Z, f, g, h)$ is a special triangle
(terminology from Remark \ref{remark-special-triangles})
as a direct summand of the distinguished hence special
triangle $(X \oplus X', Y \oplus Y', Z \oplus Z',
f \oplus f', g \oplus g', h \oplus h')$. Using TR1 let
$(X, Y, Q, f, g'', h'')$ be a distinguished triangle. By TR3 there exists
a morphism of distinguished triangles
 $(X \oplus X', Y \oplus Y', Z \oplus Z',
f \oplus f', g \oplus g', h \oplus h') \to (X, Y, Q, f, g'', h'')$.
Composing this with the inclusion map we get a morphism of triangles
$$
(1, 1, c) :
(X, Y, Z, f, g, h)
\longrightarrow
(X, Y, Q, f, g'', h'')
$$
By Remark \ref{remark-special-triangles}
we find that $c$ is an isomorphism and we conclude
that (2) holds.
\end{proof}

\begin{lemma}
\label{lemma-split}
Let $\mathcal{D}$ be a pre-triangulated category.
Let $(X, Y, Z, f, g, h)$ be a distinguished triangle.
\begin{enumerate}
\item If $h = 0$, then there exists a right inverse $s : Z \to Y$ to $g$.
\item For any right inverse $s : Z \to Y$ of $g$ the map
$f \oplus s : X \oplus Z \to Y$ is an isomorphism.
\item For any objects $X', Z'$ of $\mathcal{D}$ the triangle
$(X', X' \oplus Z', Z', (1, 0), (0, 1), 0)$ is distinguished.
\end{enumerate}
\end{lemma}

\begin{proof}
To see (1) use that
$\Hom_\mathcal{D}(Z, Y) \to \Hom_\mathcal{D}(Z, Z) \to
\Hom_\mathcal{D}(Z, X[1])$
is exact by
Lemma \ref{lemma-representable-homological}.
By the same token, if $s$ is as in (2), then $h = 0$ and the sequence
$$
0 \to \Hom_\mathcal{D}(W, X) \to \Hom_\mathcal{D}(W, Y)
\to \Hom_\mathcal{D}(W, Z) \to 0
$$
is split exact (split by $s : Z \to Y$). Hence by Yoneda's lemma we
see that $X \oplus Z \to Y$ is an isomorphism. The last assertion follows
from TR1 and
Lemma \ref{lemma-direct-sum-triangles}.
\end{proof}

\begin{lemma}
\label{lemma-when-split}
Let $\mathcal{D}$ be a pre-triangulated category.
Let $f : X \to Y$ be a morphism of $\mathcal{D}$.
The following are equivalent
\begin{enumerate}
\item $f$ has a kernel,
\item $f$ has a cokernel,
\item $f$ is the isomorphic to a composition
$K \oplus Z \to Z \to Z \oplus Q$ of a projection and coprojection
for some objects $K, Z, Q$ of $\mathcal{D}$.
\end{enumerate}
\end{lemma}

\begin{proof}
Any morphism isomorphic to a map of the form
$X' \oplus Z \to Z \oplus Y'$ has both a kernel and a cokernel.
Hence (3) $\Rightarrow$ (1), (2).
Next we prove (1) $\Rightarrow$ (3).
Suppose first that $f : X \to Y$ is a monomorphism, i.e., its kernel is zero.
By TR1 there exists a distinguished triangle $(X, Y, Z, f, g, h)$.
By Lemma \ref{lemma-composition-zero} the composition
$f \circ h[-1] = 0$. As $f$ is a monomorphism we see that $h[-1] = 0$
and hence $h = 0$. Then
Lemma \ref{lemma-split}
implies that $Y = X \oplus Z$, i.e., we see that (3) holds.
Next, assume $f$ has a kernel $K$. As $K \to X$ is a monomorphism we
conclude $X = K \oplus X'$ and $f|_{X'} : X' \to Y$ is a monomorphism.
Hence $Y = X' \oplus Y'$ and we win.
The implication (2) $\Rightarrow$ (3) is dual to this.
\end{proof}

\begin{lemma}
\label{lemma-products-sums-shifts-triangles}
Let $\mathcal{D}$ be a pre-triangulated category.
Let $I$ be a set.
\begin{enumerate}
\item Let $X_i$, $i \in I$ be a family of objects of $\mathcal{D}$.
\begin{enumerate}
\item If $\prod X_i$ exists, then $(\prod X_i)[1] = \prod X_i[1]$.
\item If $\bigoplus X_i$ exists, then $(\bigoplus X_i)[1] = \bigoplus X_i[1]$.
\end{enumerate}
\item Let $X_i \to Y_i \to Z_i \to X_i[1]$ be a family of distinguished
triangles of $\mathcal{D}$.
\begin{enumerate}
\item If $\prod X_i$, $\prod Y_i$, $\prod Z_i$ exist, then
$\prod X_i \to \prod Y_i \to \prod Z_i \to \prod X_i[1]$
is a distinguished triangle.
\item If $\bigoplus X_i$, $\bigoplus Y_i$,
$\bigoplus Z_i$ exist, then
$\bigoplus X_i \to \bigoplus Y_i \to \bigoplus Z_i \to \bigoplus X_i[1]$
is a distinguished triangle.
\end{enumerate}
\end{enumerate}
\end{lemma}

\begin{proof}
Part (1) is true because $[1]$ is an autoequivalence of $\mathcal{D}$
and because direct sums and products are defined in terms of the
category structure. Let us prove (2)(a). Choose a distinguished triangle
$\prod X_i \to \prod Y_i \to Z \to \prod X_i[1]$. For each $j$ we can
use TR3 to choose a morphism $p_j : Z \to Z_j$
fitting into a morphism of distinguished
triangles with the projection maps $\prod X_i \to X_j$ and $\prod Y_i \to Y_j$.
Using the definition of products we obtain a map
$\prod p_i : Z \to \prod Z_i$ fitting into a morphism
of triangles from the distinguished triangle to the triangle
made out of the products. Observe that the ``product'' triangle
$\prod X_i \to \prod Y_i \to \prod Z_i \to \prod X_i[1]$
is special in the terminology of Remark \ref{remark-special-triangles}
because products of exact sequences of abelian groups are exact.
Hence Remark \ref{remark-special-triangles} shows that
the morphism of triangles is an isomorphism and we conclude by TR1.
The proof of (2)(b) is dual.
\end{proof}

\begin{lemma}
\label{lemma-projectors-have-images-triangulated}
Let $\mathcal{D}$ be a pre-triangulated category.
If $\mathcal{D}$ has countable products, then $\mathcal{D}$
is Karoubian.
If $\mathcal{D}$ has countable coproducts, then $\mathcal{D}$
is Karoubian.
\end{lemma}

\begin{proof}
Assume $\mathcal{D}$ has countable products. By
Homology, Lemma \ref{homology-lemma-projectors-have-images}
it suffices to check that morphisms which have a right inverse have kernels.
Any morphism which has a right inverse is an epimorphism, hence
has a kernel by
Lemma \ref{lemma-when-split}.
The second statement is dual to the first.
\end{proof}

\noindent
The following lemma makes it slightly easier to prove that a
pre-triangulated category is triangulated.

\begin{lemma}
\label{lemma-easier-axiom-four}
Let $\mathcal{D}$ be a pre-triangulated category.
In order to prove TR4 it suffices to show that given
any pair of composable morphisms
$f : X \to Y$ and $g : Y \to Z$ there exist
\begin{enumerate}
\item isomorphisms $i : X' \to X$, $j : Y' \to Y$ and
$k : Z' \to Z$, and then setting $f' = j^{-1}fi : X' \to Y'$ and
$g' = k^{-1}gj : Y' \to Z'$ there exist
\item distinguished triangles
$(X', Y', Q_1, f', p_1, d_1)$,
$(X', Z', Q_2, g' \circ f', p_2, d_2)$
and
$(Y', Z', Q_3, g', p_3, d_3)$,
such that the assertion of TR4 holds.
\end{enumerate}
\end{lemma}

\begin{proof}
The replacement of $X, Y, Z$ by $X', Y', Z'$ is harmless by our
definition of distinguished triangles and their isomorphisms.
The lemma follows from the fact that the distinguished triangles
$(X', Y', Q_1, f', p_1, d_1)$,
$(X', Z', Q_2, g' \circ f', p_2, d_2)$
and
$(Y', Z', Q_3, g', p_3, d_3)$
are unique up to isomorphism by
Lemma \ref{lemma-cone-triangle-unique-isomorphism}.
\end{proof}

\begin{lemma}
\label{lemma-triangulated-subcategory}
Let $\mathcal{D}$ be a pre-triangulated category.
Assume that $\mathcal{D}'$ is an additive full subcategory of $\mathcal{D}$.
The following are equivalent
\begin{enumerate}
\item there exists a set of triangles $\mathcal{T}'$ such that
$(\mathcal{D}', \mathcal{T}')$ is a pre-triangulated subcategory
of $\mathcal{D}$,
\item $\mathcal{D}'$ is preserved under $[1]$ and
$[1] : \mathcal{D}' \to \mathcal{D}'$ is an auto-equivalence and
given any morphism $f : X \to Y$ in $\mathcal{D}'$ there exists
a distinguished triangle $(X, Y, Z, f, g, h)$ in $\mathcal{D}$
such that $Z$ is isomorphic to an object of $\mathcal{D}'$.
\end{enumerate}
In this case $\mathcal{T}'$ as in (1) is the set of distinguished triangles
$(X, Y, Z, f, g, h)$ of $\mathcal{D}$ such that
$X, Y, Z \in \Ob(\mathcal{D}')$. Finally, if $\mathcal{D}$
is a triangulated category, then (1) and (2) are also equivalent to
\begin{enumerate}
\item[(3)] $\mathcal{D}'$ is a triangulated subcategory.
\end{enumerate}
\end{lemma}

\begin{proof}
Omitted.
\end{proof}

\begin{lemma}
\label{lemma-exact-functor-additive}
An exact functor of pre-triangulated categories is additive.
\end{lemma}

\begin{proof}
Let $F : \mathcal{D} \to \mathcal{D}'$ be an exact functor of
pre-triangulated categories. Since
$(0, 0, 0, 1_0, 1_0, 0)$ is a distinguished triangle of $\mathcal{D}$
the triangle
$$
(F(0), F(0), F(0), 1_{F(0)}, 1_{F(0)}, F(0))
$$
is distinguished in $\mathcal{D}'$.
This implies that $1_{F(0)} \circ 1_{F(0)}$ is zero, see
Lemma \ref{lemma-composition-zero}.
Hence $F(0)$ is the zero object of $\mathcal{D}'$. This also implies
that $F$ applied to any zero morphism is zero (since a morphism in
an additive category is zero if and only if it factors through the
zero object). Next, using that
$(X, X \oplus Y, Y, (1, 0), (0, 1), 0)$ is a distinguished triangle by
Lemma \ref{lemma-split} part (3), we see that
$(F(X), F(X \oplus Y), F(Y), F(1, 0), F(0, 1), 0)$ is
one too. This implies that the map $F(X) \oplus F(Y) \to F(X \oplus Y)$
is an isomorphism by Lemma \ref{lemma-split} part (2). To finish we apply
Homology, Lemma \ref{homology-lemma-additive-functor}.
\end{proof}

\begin{lemma}
\label{lemma-exact-equivalence}
Let $F : \mathcal{D} \to \mathcal{D}'$ be a fully faithful exact functor
of pre-triangulated categories. Then a triangle $(X, Y, Z, f, g, h)$
of $\mathcal{D}$ is distinguished if and only if
$(F(X), F(Y), F(Z), F(f), F(g), F(h))$ is distinguished in $\mathcal{D}'$.
\end{lemma}

\begin{proof}
The ``only if'' part is clear. Assume $(F(X), F(Y), F(Z))$ is
distinguished in $\mathcal{D}'$. Pick a distinguished triangle
$(X, Y, Z', f, g', h')$ in $\mathcal{D}$. By
Lemma \ref{lemma-cone-triangle-unique-isomorphism}
there exists an isomorphism of triangles
$$
(1, 1, c') : (F(X), F(Y), F(Z)) \longrightarrow (F(X), F(Y), F(Z')).
$$
Since $F$ is fully faithful, there exists a morphism $c : Z \to Z'$
such that $F(c) = c'$. Then $(1, 1, c)$ is an isomorphism between
$(X, Y, Z)$ and $(X, Y, Z')$. Hence $(X, Y, Z)$ is distinguished
by TR1.
\end{proof}

\begin{lemma}
\label{lemma-composition-exact}
Let $\mathcal{D}, \mathcal{D}', \mathcal{D}''$ be pre-triangulated categories.
Let $F : \mathcal{D} \to \mathcal{D}'$ and
$F' : \mathcal{D}' \to \mathcal{D}''$ be exact functors.
Then $F' \circ F$ is an exact functor.
\end{lemma}

\begin{proof}
Omitted.
\end{proof}

\begin{lemma}
\label{lemma-exact-compose-homological-functor}
Let $\mathcal{D}$ be a pre-triangulated category.
Let $\mathcal{A}$ be an abelian category.
Let $H : \mathcal{D} \to \mathcal{A}$ be a homological functor.
\begin{enumerate}
\item Let $\mathcal{D}'$ be a pre-triangulated category.
Let $F : \mathcal{D}' \to \mathcal{D}$ be an exact functor.
Then the composition $H \circ F$ is a homological functor as well.
\item Let $\mathcal{A}'$ be an abelian category. Let
$G : \mathcal{A} \to \mathcal{A}'$ be an exact functor.
Then $G \circ H$ is a homological functor as well.
\end{enumerate}
\end{lemma}

\begin{proof}
Omitted.
\end{proof}

\begin{lemma}
\label{lemma-exact-compose-delta-functor}
Let $\mathcal{D}$ be a triangulated category.
Let $\mathcal{A}$ be an abelian category.
Let $G : \mathcal{A} \to \mathcal{D}$ be a $\delta$-functor.
\begin{enumerate}
\item Let $\mathcal{D}'$ be a triangulated category.
Let $F : \mathcal{D} \to \mathcal{D}'$ be an exact functor.
Then the composition $F \circ G$ is a $\delta$-functor as well.
\item Let $\mathcal{A}'$ be an abelian category. Let
$H : \mathcal{A}' \to \mathcal{A}$ be an exact functor.
Then $G \circ H$ is a $\delta$-functor as well.
\end{enumerate}
\end{lemma}

\begin{proof}
Omitted.
\end{proof}

\begin{lemma}
\label{lemma-compose-delta-functor-homological}
Let $\mathcal{D}$ be a triangulated category.
Let $\mathcal{A}$ and $\mathcal{B}$ be abelian categories.
Let $G : \mathcal{A} \to \mathcal{D}$ be a $\delta$-functor.
Let $H : \mathcal{D} \to \mathcal{B}$ be a homological functor.
Assume that $H^{-1}(G(A)) = 0$ for all $A$ in $\mathcal{A}$.
Then the collection
$$
\{H^n \circ G, H^n(\delta_{A \to B \to C})\}_{n \geq 0}
$$
is a $\delta$-functor from $\mathcal{A} \to \mathcal{B}$, see
Homology, Definition \ref{homology-definition-cohomological-delta-functor}.
\end{lemma}

\begin{proof}
The notation signifies the following. If
$0 \to A \xrightarrow{a} B \xrightarrow{b} C \to 0$ is
a short exact sequence in $\mathcal{A}$, then
$$
\delta = \delta_{A \to B \to C} : G(C) \to G(A)[1]
$$
is a morphism in $\mathcal{D}$ such that
$(G(A), G(B), G(C), a, b, \delta)$ is
a distinguished triangle, see
Definition \ref{definition-delta-functor}.
Then $H^n(\delta) : H^n(G(C)) \to H^n(G(A)[1]) = H^{n + 1}(G(A))$
is clearly functorial in the short exact sequence.
Finally, the long exact cohomology sequence
(\ref{equation-long-exact-cohomology-sequence})
combined with the vanishing of $H^{-1}(G(C))$
gives a long exact sequence
$$
0 \to H^0(G(A)) \to H^0(G(B)) \to H^0(G(C))
\xrightarrow{H^0(\delta)} H^1(G(A)) \to \ldots
$$
in $\mathcal{B}$ as desired.
\end{proof}

\noindent
The proof of the following result uses TR4.

\begin{proposition}
\label{proposition-9}
Let $\mathcal{D}$ be a triangulated category. Any commutative diagram
$$
\xymatrix{
X \ar[r] \ar[d] & Y \ar[d] \\
X' \ar[r] & Y'
}
$$
can be extended to a diagram
$$
\xymatrix{
X \ar[r] \ar[d] & Y \ar[r] \ar[d] & Z \ar[r] \ar[d] & X[1] \ar[d] \\
X' \ar[r] \ar[d] & Y' \ar[r] \ar[d] & Z' \ar[r] \ar[d] & X'[1] \ar[d] \\
X'' \ar[r] \ar[d] & Y'' \ar[r] \ar[d] & Z'' \ar[r] \ar[d] & X''[1] \ar[d] \\
X[1] \ar[r] & Y[1] \ar[r] & Z[1] \ar[r] & X[2]
}
$$
where all the squares are commutative, except for the lower right square
which is anticommutative. Moreover, each of the rows and columns are
distinguished triangles. Finally, the morphisms on the bottom row
(resp.\ right column) are obtained from the morphisms of the top row
(resp.\ left column) by applying $[1]$.
\end{proposition}

\begin{proof}
During this proof we avoid writing the arrows in order to make the proof
legible. Choose distinguished triangles
$(X, Y, Z)$, $(X', Y', Z')$, $(X, X', X'')$, $(Y, Y', Y'')$, and
$(X, Y', A)$. Note that the morphism $X \to Y'$ is both equal
to the composition $X \to Y \to Y'$ and equal to the composition
$X \to X' \to Y'$. Hence, we can find morphisms
\begin{enumerate}
\item $a : Z \to A$ and $b : A \to Y''$, and
\item $a' : X'' \to A$ and $b' : A \to Z'$
\end{enumerate}
as in TR4. Denote $c : Y'' \to Z[1]$ the composition
$Y'' \to Y[1] \to Z[1]$ and denote $c' : Z' \to X''[1]$ the composition
$Z' \to X'[1] \to X''[1]$. The conclusion of our application TR4
are that
\begin{enumerate}
\item $(Z, A, Y'', a, b, c)$, $(X'', A, Z', a', b', c')$
are distinguished triangles,
\item $(X, Y, Z) \to (X, Y', A)$,
$(X, Y', A) \to (Y, Y', Y'')$,
$(X, X', X'') \to (X, Y', A)$,
$(X, Y', A) \to (X', Y', Z')$
are morphisms of triangles.
\end{enumerate}
First using that
$(X, X', X'') \to (X, Y', A)$ and $(X, Y', A) \to (Y, Y', Y'')$.
are morphisms of triangles we see the first of the diagrams
$$
\vcenter{
\xymatrix{
X' \ar[r] \ar[d] & Y' \ar[d] \\
X'' \ar[r]^{b \circ a'} \ar[d] & Y'' \ar[d] \\
X[1] \ar[r] & Y[1]
}
}
\quad\text{and}\quad
\vcenter{
\xymatrix{
Y \ar[r] \ar[d] & Z \ar[d]^{b' \circ a} \ar[r] & X[1] \ar[d] \\
Y' \ar[r] & Z' \ar[r] & X'[1]
}
}
$$
is commutative. The second is commutative too using that
$(X, Y, Z) \to (X, Y', A)$ and $(X, Y', A) \to (X', Y', Z')$ are morphisms
of triangles. At this point we choose a distinguished triangle
$(X'', Y'' , Z'')$ starting with the map $b \circ a' : X'' \to Y''$.

\medskip\noindent
Next we apply TR4 one more time to the morphisms
$X'' \to A \to Y''$ and the triangles
$(X'', A, Z', a', b', c')$,
$(X'', Y'', Z'')$, and
$(A, Y'', Z[1], b, c , -a[1])$ to get morphisms
$a'' : Z' \to Z''$ and $b'' : Z'' \to Z[1]$.
Then $(Z', Z'', Z[1], a'', b'', - b'[1] \circ a[1])$ is a distinguished
triangle, hence also $(Z, Z', Z'', -b' \circ a, a'', -b'')$
and hence also $(Z, Z', Z'', b' \circ a, a'', b'')$.
Moreover, $(X'', A, Z') \to (X'', Y'', Z'')$ and
$(X'', Y'', Z'') \to (A, Y'', Z[1], b, c , -a[1])$
are morphisms of triangles.
At this point we have defined all the distinguished triangles
and all the morphisms, and all that's left is to verify some
commutativity relations.

\medskip\noindent
To see that the middle square in the diagram commutes, note
that the arrow $Y' \to Z'$ factors as $Y' \to A \to Z'$
because $(X, Y', A) \to (X', Y', Z')$ is a morphism of triangles.
Similarly, the morphism $Y' \to Y''$ factors as
$Y' \to A \to Y''$ because $(X, Y', A) \to (Y, Y', Y'')$ is a
morphism of triangles. Hence the middle square commutes because
the square with sides $(A, Z', Z'', Y'')$ commutes as
$(X'', A, Z') \to (X'', Y'', Z'')$ is a morphism of triangles (by TR4).
The square with sides $(Y'', Z'', Y[1], Z[1])$ commutes
because $(X'', Y'', Z'') \to (A, Y'', Z[1], b, c , -a[1])$
is a morphism of triangles and $c : Y'' \to Z[1]$ is the composition
$Y'' \to Y[1] \to Z[1]$.
The square with sides $(Z', X'[1], X''[1], Z'')$ is commutative
because $(X'', A, Z') \to (X'', Y'', Z'')$ is a morphism of triangles
and $c' : Z' \to X''[1]$ is the composition $Z' \to X'[1] \to X''[1]$.
Finally, we have to show that the square with sides
$(Z'', X''[1], Z[1], X[2])$ anticommutes. This holds because
$(X'', Y'', Z'') \to (A, Y'', Z[1], b, c , -a[1])$
is a morphism of triangles and we're done.
\end{proof}








